\chapter{Conclusões e Trabalho Futuro} \label{chap:concl}

\section*{}

Deve ser apresentado um resumo do trabalho realizado e apreciada a
satisfação dos objetivos do trabalho, uma lista de contribuições
principais do trabalho e as direções para trabalho futuro.

A escrita deste capítulo deve ser orientada para a total compreensão
do trabalho, tendo em atenção que, depois de ler o Resumo e a
Introdução, a maioria dos leitores passará à leitura deste capítulo de
conclusões e recomendações para trabalho futuro.

\section{Achievements}

Concretizações:
	Provar a necessidade de explorar visualizações que incorporam partituras alinhadas com tempo real de uma performance/gravação.
	Provar o valor de usar uma ferramenta modular em python.

\section{Future Work}

Problemas da solução:
	O problema do sistema estar dependente do SVG no formato em que sai do ScoreWarp.
	Dificuldade em obter os ficheiros correctos (MAPS, MEI, etc)
	O problema de nao conseguir definir o espaçamento da partitura.
	O problema da representação do espetrograma como PNG.
Trabalho Futuro:
	Integração com um GUI. 
	Problema da paginação
	Integração para visualização de outros parametros MIR.
	Integração de metodos para gerar visualizações costumizaveis sem ter de alterar manualmente o SVG.


\vspace*{12mm}
